% 06_tablas/tabla_cdiu.tex
% Tabla 1. Cuadro de Categorías, Dimensiones, Instrumentos y Unidades de Análisis (CDIU)
% Versión simplificada para el Capítulo I — nueve categorías principales
% Usa longtable para permitir que la tabla se extienda en varias páginas

\begingroup
\small
\begin{longtable}{|p{0.03\textwidth}|p{0.17\textwidth}|p{0.17\textwidth}|p{0.21\textwidth}|p{0.23\textwidth}|}

\caption{Cuadro de Categorías, Dimensiones, Instrumentos y Unidades de Análisis (CDIU)}
\label{tab:cdiu} \\

\hline
\textbf{N.°} & \textbf{Categoría} & \textbf{Dimensión principal} & \textbf{Instrumento de análisis} & \textbf{Unidad de análisis} \\
\hline
\endfirsthead

\multicolumn{5}{l}{\small\textit{(continuación de la Tabla \ref*{tab:cdiu})}} \\[2pt]
\hline
\textbf{N.°} & \textbf{Categoría} & \textbf{Dimensión principal} & \textbf{Instrumento de análisis} & \textbf{Unidad de análisis} \\
\hline
\endhead

\hline
\endfoot

1 &
Derecho ambiental ecuatoriano &
Marco constitucional, legal y jurisprudencial &
Análisis normativo constitucional y legislativo; análisis jurisprudencial &
Constitución de la República del Ecuador (2008), arts. 14, 15, 395--399, 407--409; Código Orgánico del Ambiente (COA); sentencias de la Corte Constitucional del Ecuador \\
\hline

2 &
Derechos de la naturaleza &
Reconocimiento constitucional y jurisprudencial &
Fichas jurisprudenciales; análisis normativo constitucional &
Arts. 71--74 CRE; Sentencia No. 1149-19-JP/21 (caso Bosque Protector Los Cedros); Sentencia No. 22-18-IN/21; Guía jurisprudencial CEDEC 2023 \\
\hline

3 &
Actividad minera privada &
Normativa, institucional y sectorial &
Matriz de análisis normativo legislativo; análisis documental institucional &
Ley de Minería (texto actualizado hasta agosto de 2025); Ley Orgánica para el Fortalecimiento de los Sectores Estratégicos de Minería y Energía (2026); documentos del Ministerio de Energía y Minas; documentos de ARCOM \\
\hline

4 &
Control estatal ambiental &
Institucional y procedimental &
Matriz de análisis normativo; análisis institucional &
Arts. 8, 9 Ley de Minería; arts. 162, 166, 174 COA; arts. 431--435 Reglamento al COA; Informe ARCOM Rendición de Cuentas 2024 \\
\hline

5 &
Seguridad jurídica &
Normativa y procedimental &
Análisis normativo constitucional; análisis doctrinal &
Art. 82 CRE; Reglamento General a la Ley de Minería (arts. 9--12); Ley de Fortalecimiento 2026; Villacís Calvas (2022) \\
\hline

6 &
Reformas legislativas y normativas mineras &
Técnica legislativa y análisis comparado &
Matriz de análisis documental legislativo; análisis normativo comparado &
Ley Orgánica Reformatoria a la Ley de Minería (2013); Ley de Fortalecimiento 2026 (R.O. Quinto Suplemento No. 234, 2 de marzo de 2026); documentos del proceso legislativo disponibles \\
\hline

7 &
Consulta ambiental y participación ciudadana &
Constitucional, jurisprudencial e internacional &
Análisis jurisprudencial; análisis de instrumentos internacionales &
Art. 398 CRE; Sentencia No. 1149-19-JP/21; Sentencia No. 22-18-IN/21; arts. 5, 6, 7 Acuerdo de Escazú; CEPAL Ruta Escazú Ecuador 2023 \\
\hline

8 &
Responsabilidad ambiental y reparación integral &
Constitucional, legal y jurisprudencial &
Análisis normativo y jurisprudencial; fichas doctrinales &
Arts. 396, 397 CRE; arts. 10, 11, 299--325 COA; arts. 257--261 COIP; Sentencia No. 1149-19-JP/21 (parte resolutiva); Medina y García (2026) \\
\hline

9 &
Propuesta de compatibilidad constitucional y ambiental &
Técnica jurídica e instrumental &
Matriz de Verificación de Compatibilidad Constitucional y Ambiental &
Criterios derivados de arts. 71--74, 82, 395--399 CRE; jurisprudencia constitucional vinculante; Acuerdo de Escazú; COA; Ley de Fortalecimiento 2026 \\
\hline

\multicolumn{5}{p{0.9\textwidth}}{\small\textit{Nota.} Elaboración propia con base en los objetivos específicos de la investigación, el corpus normativo, jurisprudencial, doctrinal e institucional del repositorio y la metodología documental, descriptiva y jurídico-crítica adoptada. La versión extendida de esta matriz, con 35 dimensiones desarrolladas, indicadores y fuentes verificables, funciona como soporte metodológico interno del Capítulo III. Para efectos del presente capítulo se presenta la versión simplificada por categorías.} \\

\end{longtable}
\endgroup
